% ======================================================================
% Diagrama de flujo — generado automáticamente por generar_diagrama_flujo.py
% NO EDITAR MANUALMENTE; regenerar desde el descriptor JSON.
% ======================================================================
\documentclass[11pt,a4paper]{article}

% ── Codificación, fuentes y lenguaje ───────────────────────────────────────────
\usepackage[utf8]{inputenc}
\usepackage[T1]{fontenc}
\usepackage[default,scale=0.95]{sourcesanspro}
\renewcommand{\familydefault}{\sfdefault}
\usepackage[spanish,es-noshorthands,es-tabla]{babel}

% ── Geometría y espaciado ─────────────────────────────────────────────────────
\usepackage[top=2.2cm, bottom=2.5cm, left=2.2cm, right=2.2cm, headheight=15pt]{geometry}
\usepackage{setspace}
\setstretch{1.18}
\usepackage{parskip}

% ── TikZ (simbolos ISO 5807 / ANSI X3.5) ──────────────────────────────────────
\usepackage{tikz}
\usetikzlibrary{babel, arrows.meta, positioning, shapes.geometric,
                shapes.symbols, decorations.pathmorphing, calc}

% ── Tablas y listas ───────────────────────────────────────────────────────────
\usepackage{booktabs}
\usepackage{tabularx}
\usepackage{array}
\usepackage{enumitem}

% ── Colores, cajas e iconos ───────────────────────────────────────────────────
\usepackage[dvipsnames]{xcolor}
\usepackage{tcolorbox}
\tcbuselibrary{skins, breakable}
\usepackage{fontawesome5}

% ── Secciones con barra de color (infografía) ─────────────────────────────────
\usepackage{titlesec}
\titleformat{\section}
  {\normalfont\LARGE\bfseries\color{azulNoche}}
  {\colorbox{cyanNeon}{\color{fondoOscuro}\thesection}}{0.7em}{}
  [\vspace{2pt}\color{cyanNeon}\rule{\linewidth}{1.8pt}]
\titlespacing*{\section}{0pt}{24pt}{10pt}
\titleformat{\subsection}
  {\normalfont\large\bfseries\color{azulNoche}}
  {\textcolor{cyanNeon}{\thesubsection}}{0.7em}{}
  [\vspace{1pt}\color{grisLinea}\rule{\linewidth}{0.6pt}]
\titlespacing*{\subsection}{0pt}{14pt}{6pt}

% ── Cabeceras y pies de página ────────────────────────────────────────────────
\usepackage{fancyhdr}
\usepackage{lastpage}
\pagestyle{fancy}
\fancyhf{}
\renewcommand{\headrulewidth}{0pt}
\renewcommand{\footrulewidth}{0pt}
\fancyfoot[C]{%
  \begin{minipage}{\linewidth}
    \vspace{2pt}
    \color{cyanNeon}\rule{\linewidth}{1.2pt}\\[2pt]
    \centering\color{grisTexto}\footnotesize
    Página \thepage\ de \pageref{LastPage}
  \end{minipage}%
}

% ── Hiperenlaces ──────────────────────────────────────────────────────────────
\usepackage[colorlinks=true, linkcolor=azulNoche, urlcolor=cyanNeon]{hyperref}
\sloppy
\emergencystretch 3em

% ── Paleta de colores ─────────────────────────────────────────────────────────
\definecolor{fondoOscuro}{HTML}{0D1B2A}
\definecolor{verdeTurquesa}{HTML}{004D40}
\definecolor{azulNoche}{HTML}{1B3A6B}
\definecolor{azulMedio}{HTML}{2A5298}
\definecolor{cyanNeon}{HTML}{00C8E0}
\definecolor{verdeSignal}{HTML}{00B887}
\definecolor{naranjaVivo}{HTML}{FF7043}
\definecolor{rojoAlerta}{HTML}{E53935}
\definecolor{grisPapel}{HTML}{F4F6FA}
\definecolor{grisLinea}{HTML}{C8D0E0}
\definecolor{grisTexto}{HTML}{6B7A99}
\definecolor{amarilloNota}{HTML}{FFB300}

% ── Tarjetas ──────────────────────────────────────────────────────────────────
\tcbset{
  cajaContenido/.style={
    enhanced, breakable,
    arc=5pt, outer arc=5pt,
    colback=grisPapel, colframe=grisLinea,
    boxrule=0.8pt,
    fonttitle=\bfseries\small, coltitle=azulNoche,
    top=6pt, bottom=6pt, left=10pt, right=10pt,
  },
}

% ── Macros ────────────────────────────────────────────────────────────────────
\newcommand{\iconotexto}[2]{\textcolor{cyanNeon}{\faIcon{#1}}~#2}
\newcommand{\iconoBanda}{sitemap}
\newcommand{\bandaTitulo}[3][fondoOscuro]{%
  \noindent
  \begin{tcolorbox}[%
    enhanced, arc=4mm, colback=#1!10!white, colframe=#1, boxrule=0pt,
    borderline west={6pt}{0pt}{cyanNeon},
    top=6pt, bottom=6pt, left=8pt, right=8pt,
  ]
    \noindent
    \begin{minipage}[c]{0.10\linewidth}
      \centering\color{cyanNeon}\faIcon{project-diagram}
    \end{minipage}%
    \hspace{8pt}%
    \begin{minipage}[c]{0.88\linewidth}
      {\fontsize{20}{24}\selectfont\bfseries\color{white}#2}\par\vspace{5pt}%
      \textcolor{cyanNeon!80!white}{\small\faIcon{angle-right}~#3}%
    \end{minipage}
    \vspace{4pt}
  \end{tcolorbox}%
  \vspace{8pt}%
}


  % ── Estilos de flujo (ISO 5807 / ANSI X3.5) ─────────────────────────────
  \tikzset{
    flujo/.style={
      >=Stealth,
      font=\footnotesize\sffamily,
    },
    terminador/.style={
      draw=azulNoche, fill=azulNoche!8!white, rounded corners=3pt,
      text width=1.6cm, align=center, inner sep=2mm,
      font=\footnotesize\sffamily,
    },
    proceso/.style={
      draw=azulNoche, fill=cyanNeon!10!white,
      text width=1.6cm, align=center, inner sep=2mm,
      font=\footnotesize\sffamily,
    },
    entradasalida/.style={
      draw=azulNoche, fill=verdeSignal!10!white,
      trapezium, trapezium left angle=75, trapezium right angle=105,
      text width=1.8cm, align=center, inner sep=2mm,
      font=\footnotesize\sffamily,
    },
    decision/.style={
      draw=azulNoche, fill=amarilloNota!20!white,
      diamond, aspect=1.6, text width=1.5cm, align=center, inner sep=1pt,
      font=\footnotesize\sffamily,
    },
    almacenamiento/.style={
      draw=azulNoche, fill=grisPapel,
      cylinder, shape border rotate=90, aspect=0.3,
      text width=1.7cm, align=center, inner sep=2mm,
      font=\footnotesize\sffamily,
    },
    documento/.style={
      draw=azulNoche, fill=azulMedio!8!white,
      text width=1.8cm, align=center, inner sep=2mm,
      font=\footnotesize\sffamily,
    },
    nota/.style={
      draw=grisLinea, dashed, fill=grisPapel,
      text width=1.8cm, align=center, inner sep=2mm,
      font=\scriptsize\sffamily,
    },
    etiqueta/.style={
      font=\scriptsize\sffamily, fill=white,
      fill opacity=0.75, text opacity=1, inner sep=1pt,
    },
  }


% ── Cabecera del documento ────────────────────────────────────────────────────
\fancyhead[L]{\small\color{azulNoche}\textbf{ELECTRÓNICA}}
\fancyhead[R]{\small\color{grisTexto}Dibujar figuras TikZ}
\renewcommand{\headrulewidth}{0pt}

% ─────────────────────────────────────────────────────────────────────────────
\begin{document}

\bandaTitulo{Figuras TikZ desde el OCR}{34 dibujos validados por compilación e inyectados en las prácticas P1-P8 de Laboratorio II}

\vspace{4pt}
\begin{center}
\begin{tcolorbox}[cajaContenido, title={\faIcon{map-signs}~Leyenda de símbolos---ISO 5807 / ANSI X3.5}]
\small
Óvalo = \textbf{terminador} (inicio/fin) $\cdot$
Rectángulo = \textbf{proceso} (acción determinista) $\cdot$
Rombo = \textbf{decisión} (ramas Sí/No) $\cdot$
Paralelogramo = \textbf{entrada/salida} (lectura/escritura de archivos) $\cdot$
Cilindro = \textbf{almacenamiento online} (estado, snapshots) $\cdot$
Rectángulo con base ondulada = \textbf{documento} (sesiones, logs) $\cdot$
Rectángulo punteado gris = \textbf{nota explicativa} (sin acción).
Flujo: arriba$\rightarrow$abajo, izquierda$\rightarrow$derecha.

\faIcon{tags}~Cada símbolo lleva una \textbf{etiqueta} (T/P/D/E/A/N + número, según su tipo);
su significado se expande en la tabla \emph{Leyenda de etiquetas} bajo cada diagrama.
\end{tcolorbox}
\end{center}

\vspace{6pt}

\section{\iconotexto{cogs}{Sección 1 - Entrada y decisión del modelo}}

\begin{center}
\begin{tikzpicture}[flujo, >=Stealth, x=1cm, y=1cm]
  % Fondo decorativo
  \fill[azulNoche!3!white, rounded corners=4pt]
    (-2.3,1.8) rectangle (5.8,-13.3);
  \node[terminador] (T1) at (0.0,0.0) {T1};
  \node[proceso] (P1) at (0.0,-2.2) {P1};
  \node[decision] (D1) at (0.0,-4.4) {D1};
  \node[almacenamiento] (A1) at (3.5,-4.4) {A1};
  \node[terminador] (T2) at (3.5,-6.6) {T2};
  \node[proceso] (P2) at (0.0,-6.6) {P2};
  \node[proceso] (P3) at (0.0,-8.8) {P3};
  \node[proceso] (P4) at (0.0,-11.0) {P4};

  \draw[->] (T1.south) -- (P1.north) node[etiqueta, left, pos=0.5]{Siempre};
  \draw[->] (P1.south) -- (D1.north) node[etiqueta, left, pos=0.5]{Siempre};
  \draw[->] (D1.east) -- (A1.west) node[etiqueta, above, pos=0.5]{Sí};
  \draw[->] (A1.south) -- (T2.north) node[etiqueta, left, pos=0.5]{Siempre};
  \draw[->] (D1.south) -- (P2.north) node[etiqueta, left, pos=0.5]{No};
  \draw[->] (P2.south) -- (P3.north) node[etiqueta, left, pos=0.5]{Siempre};
  \draw[->] (P3.south) -- (P4.north) node[etiqueta, left, pos=0.5]{Siempre};
\end{tikzpicture}
\end{center}

\vspace{6pt}
\begin{tcolorbox}[cajaContenido, title={\faIcon{table}~Leyenda de etiquetas}]
\small
\begin{tabularx}{\linewidth}{@{}>{\centering\arraybackslash}p{1.4cm}X@{}}
\toprule
\textbf{Etiqueta} & \textbf{Significado} \\ \midrule
A1 & Manifest cacheado en .tmp/figuras\_tikz/figuras.json: se reutiliza sin coste. \\
D1 & ¿Existe el manifest completo y no se pidió --forzar? \\
P1 & Validar el OCR sin API: 34 marcadores [FIGURA] repartidos en 8 prácticas. \\
P2 & Enrutador determinista: task=diseño, OCR medido, criticidad; devuelve opus. \\
P3 & Extraer descripciones con la numeración canónica pNN\_fNN, sin redefinirla. \\
P4 & Armar un lote por práctica con la lista \{id: descripción\} y el prompt acotado. \\
T1 & Inicio del flujo: markdown del OCR y directorios de manifest y prácticas. \\
T2 & Salida anticipada: cero llamadas a OpenRouter, saltar a la sección 3. \\
\bottomrule
\end{tabularx}
\end{tcolorbox}

\section{\iconotexto{draw-polygon}{Sección 2 - Generación validada por compilación}}

\begin{center}
\begin{tikzpicture}[flujo, >=Stealth, x=1cm, y=1cm]
  % Fondo decorativo
  \fill[azulNoche!3!white, rounded corners=4pt]
    (-2.3,1.8) rectangle (9.3,-15.5);
  \node[terminador] (T3) at (0.0,0.0) {T3};
  \node[proceso] (P5) at (0.0,-2.2) {P5};
  \node[decision] (D2) at (0.0,-4.4) {D2};
  \node[proceso] (P6) at (3.5,-4.4) {P6};
  \node[terminador] (T4) at (7.0,-4.4) {T4};
  \node[proceso] (P7) at (0.0,-6.6) {P7};
  \node[decision] (D3) at (0.0,-8.8) {D3};
  \node[proceso] (P8) at (0.0,-11.0) {P8};
  \node[decision] (D4) at (0.0,-13.2) {D4};
  \node[almacenamiento] (A2) at (3.5,-13.2) {A2};

  \draw[->] (T3.south) -- (P5.north) node[etiqueta, left, pos=0.5]{Siempre};
  \draw[->] (P5.south) -- (D2.north) node[etiqueta, left, pos=0.5]{Siempre};
  \draw[->] (D2.east) -- (P6.west) node[etiqueta, above, pos=0.5]{No};
  \draw[->] (P6.west) -- (D2.east);
  \draw[->] (P6.east) -- (T4.west) node[etiqueta, above, pos=0.5]{3 intentos};
  \draw[->] (D2.south) -- (P7.north) node[etiqueta, left, pos=0.5]{Sí};
  \draw[->] (P7.south) -- (D3.north) node[etiqueta, left, pos=0.5]{Siempre};
  \draw[->] (D3.east) -- (P6.west) node[etiqueta, above, pos=0.5]{No};
  \draw[->] (D3.south) -- (P8.north) node[etiqueta, left, pos=0.5]{Sí};
  \draw[->] (P8.south) -- (D4.north) node[etiqueta, left, pos=0.5]{Siempre};
  \draw[->] (D4.east) -- (P6.west) node[etiqueta, above, pos=0.5]{No};
  \draw[->] (D4.east) -- (A2.west) node[etiqueta, above, pos=0.5]{Sí};
\end{tikzpicture}
\end{center}

\vspace{6pt}
\begin{tcolorbox}[cajaContenido, title={\faIcon{table}~Leyenda de etiquetas}]
\small
\begin{tabularx}{\linewidth}{@{}>{\centering\arraybackslash}p{1.4cm}X@{}}
\toprule
\textbf{Etiqueta} & \textbf{Significado} \\ \midrule
A2 & Escribir de forma atómica código, sha y modelo en el manifest por figura. \\
D2 & ¿Bloques <<<FIGURA:id>>> completos, sin IDs ausentes, repetidos ni extraños? \\
D3 & ¿Pasa la validación determinista? \\
D4 & ¿Compila limpio el bloque? \\
P5 & Llamar a OpenRouter con el modelo decidido y max\_tokens del .env. \\
P6 & Reintentar el lote con el detalle de los fallos; presupuesto máximo 3 intentos. \\
P7 & Validación determinista: sin pgfplots, galvanometer, pC ni entornos desbalanceados. \\
P8 & Compilar la figura aislada con el preámbulo del repo; exigir 0 '!' y PDF. \\
T3 & Entrada a la generación: un lote de figuras pendientes. \\
T4 & Agotado el reintento: manifest status=error; nada se publica. \\
\bottomrule
\end{tabularx}
\end{tcolorbox}

\section{\iconotexto{flag-checkered}{Sección 3 - Inyección y entrega}}

\begin{center}
\begin{tikzpicture}[flujo, >=Stealth, x=1cm, y=1cm]
  % Fondo decorativo
  \fill[azulNoche!3!white, rounded corners=4pt]
    (-2.3,1.8) rectangle (5.8,-15.5);
  \node[terminador] (T5) at (0.0,0.0) {T5};
  \node[proceso] (P9) at (0.0,-2.2) {P9};
  \node[decision] (D5) at (0.0,-4.4) {D5};
  \node[terminador] (T6) at (3.5,-4.4) {T6};
  \node[proceso] (P10) at (0.0,-6.6) {P10};
  \node[proceso] (P11) at (0.0,-8.8) {P11};
  \node[decision] (D6) at (0.0,-11.0) {D6};
  \node[terminador] (T7) at (3.5,-11.0) {T7};
  \node[almacenamiento] (A3) at (0.0,-13.2) {A3};

  \draw[->] (T5.south) -- (P9.north) node[etiqueta, left, pos=0.5]{Siempre};
  \draw[->] (P9.south) -- (D5.north) node[etiqueta, left, pos=0.5]{Siempre};
  \draw[->] (D5.east) -- (T6.west) node[etiqueta, above, pos=0.5]{No};
  \draw[->] (D5.south) -- (P10.north) node[etiqueta, left, pos=0.5]{Sí};
  \draw[->] (P10.south) -- (P11.north) node[etiqueta, left, pos=0.5]{Siempre};
  \draw[->] (P11.south) -- (D6.north) node[etiqueta, left, pos=0.5]{Siempre};
  \draw[->] (D6.east) -- (T7.west) node[etiqueta, above, pos=0.5]{No};
  \draw[->] (D6.south) -- (A3.north) node[etiqueta, left, pos=0.5]{Sí};
\end{tikzpicture}
\end{center}

\vspace{6pt}
\begin{tcolorbox}[cajaContenido, title={\faIcon{table}~Leyenda de etiquetas}]
\small
\begin{tabularx}{\linewidth}{@{}>{\centering\arraybackslash}p{1.4cm}X@{}}
\toprule
\textbf{Etiqueta} & \textbf{Significado} \\ \midrule
A3 & Artefactos finales .tex/.pdf, sha\_conjunto y vista por corrida; alerta de éxito. \\
D5 & ¿Están todas las figuras del manifest? \\
D6 & ¿Entrega completa y limpia? \\
P10 & Regenerar P1-P8 con --figuras: inyectar el bloque donde estaba el marcador. \\
P11 & Compilar las 8 prácticas y verificar el log: 0 líneas de error y 0 avisos Overfull. \\
P9 & Validar cobertura: status ok y un código validado por cada una de las 34 figuras. \\
T5 & Entrada a la entrega con el manifest ya validado. \\
T6 & Falta alguna: se cae a la caja de descripción, nunca se publica un dibujo sin validar. \\
T7 & Fallo categorizado: se registra en la vista run\_state y se alerta. \\
\bottomrule
\end{tabularx}
\end{tcolorbox}


\section{\iconotexto{sticky-note}{Notas}}
\begin{itemize}
\item El modelo lo decide execution/enrutador.py con --task 'diseño' --archivos <OCR> --critico: mismo descriptor, mismo tier, siempre.
\item La numeración pNN\_fNN se importa de generar\_practicas\_desde\_ocr.py: generador e inyector recorren exactamente el mismo orden.
\item Ningún dibujo entra al manifest sin compilar aislado con 0 líneas '!' en el .log; con el manifest presente no se llama a la API.
\end{itemize}

% ─────────────────────────────────────────────────────────────────────────────
\end{document}
